\definecolor{caseborder}{HTML}{D6DEE6}
\definecolor{casefill}{HTML}{F8FAFB}
\definecolor{caseheadfill}{HTML}{EEF3F6}
\definecolor{caseblue}{HTML}{2F5F83}
\definecolor{caseteal}{HTML}{2F7E72}
\definecolor{caseconstruct}{HTML}{9B4F3A}
\definecolor{caseactive}{HTML}{8A6A1F}
\definecolor{caseref}{HTML}{65707A}
\definecolor{casepale}{HTML}{E8EEF2}

\newcommand{\casebadge}[3]{%
  \begingroup
  \setlength{\fboxsep}{1.8pt}%
  \colorbox{#1}{\textcolor{#2}{\scriptsize\strut\textbf{#3}}}%
  \endgroup}
\newcommand{\casecode}[1]{\casebadge{casepale}{caseblue}{#1}}
\newcommand{\userturn}[1]{\casebadge{caseblue}{white}{\(\triangleright\) #1}\par{\scriptsize\textcolor{caseblue}{User}}}
\newcommand{\assistantturn}[1]{\casebadge{caseteal}{white}{\(\diamond\) #1}\par{\scriptsize\textcolor{caseteal}{Assistant}}}
\newcommand{\activebadge}{\casebadge{caseactive}{white}{\(\rightarrow\) A active}}
\newcommand{\constructivebadge}{\casebadge{caseconstruct}{white}{\(\uparrow\) C constructive}}
\newcommand{\scaffoldbadge}{\casebadge{caseteal}{white}{\(\ast\) Scaffolded}}
\newcommand{\referencebadge}{\casebadge{caseref}{white}{\(\cdot\) Reference}}
\newcommand{\casekey}[1]{\textcolor{caseconstruct}{\textbf{#1}}}
\newcommand{\supportkey}[1]{\textcolor{caseteal}{\textbf{#1}}}
\newcommand{\caselegend}{%
  \begingroup
  {\scriptsize\textcolor{caseblue}{\textbf{Encoding}}\quad
  \constructivebadge\hspace{3pt}\activebadge\hspace{3pt}\scaffoldbadge\hspace{3pt}\referencebadge}\par\vspace{2pt}
  {\scriptsize\textcolor{caseblue}{\textbf{Codes}}\quad
  \casecode{I2} support intent; \casecode{M} support-form labels may co-occur.}
  \endgroup}

\begin{table*}[p]
\footnotesize
\centering
\caption{\textbf{Complete de-identified coding case with turn-level labels.} The illustrative transcript is derived from a retained conversation and lightly edited while preserving turn sequence, task structure and annotation decisions. Colored badges encode role, user engagement level and assistant support category.}
\label{tab:complete_coding_case}
\setlength{\fboxsep}{5pt}
\setlength{\fboxrule}{0.6pt}
\begin{adjustbox}{center,max width=0.98\textwidth}
\fcolorbox{caseborder}{casefill}{%
\begin{minipage}{0.955\textwidth}
\setlength{\tabcolsep}{3.5pt}
\renewcommand{\arraystretch}{1.18}
\begingroup
\setlength{\fboxsep}{3pt}
\colorbox{caseheadfill}{%
\begin{minipage}{0.975\linewidth}
\textcolor{caseblue}{\textbf{Coding case}}\quad \textbf{SQL transaction handling}\\[-1pt]
{\scriptsize WildChat coding; intentional framing; four turns. Constructive user follow-up probes a programming mechanism; scaffolded support provides explanation, procedure and modelling.}
\end{minipage}}
\endgroup
\vspace{5pt}

\caselegend
\vspace{5pt}

\begin{tabularx}{\linewidth}{>{\raggedright\arraybackslash}p{0.115\linewidth}>{\raggedright\arraybackslash}X>{\raggedright\arraybackslash}p{0.305\linewidth}}
\toprule
\textbf{Turn} & \textbf{De-identified transcript} & \textbf{Label and annotation cue} \\
\midrule
\userturn{U1} & ``How does \texttt{XACT\_ABORT} work in a SQL Server script? If my delete statements are inside a transaction so failed deletes should roll back, should \texttt{XACT\_ABORT} be on or off? Does it apply only to the transaction or to the whole script?'' & \constructivebadge\par The user probes a \casekey{mechanism}, asks about \casekey{scope} and links the answer to rollback behaviour. \\
\addlinespace[5pt]
\assistantturn{A1} & The assistant explains that \texttt{XACT\_ABORT ON} makes many run-time errors abort and roll back the current transaction, distinguishes transaction-level rollback from statements outside the transaction and explains why the setting is commonly used for update/delete scripts. & \scaffoldbadge\quad \casecode{I2}\ \casecode{M4}\par The reply \supportkey{explains the mechanism} connecting errors, transactions and rollback. \\
\addlinespace[5pt]
\userturn{U2} & ``Do I need to set \texttt{XACT\_ABORT OFF} at the end because of scope? What about using \texttt{TRY}/\texttt{CATCH} outside a transaction instead?'' & \constructivebadge\par The user extends the prior answer by comparing \casekey{scope} and an \casekey{alternative strategy}. \\
\addlinespace[5pt]
\assistantturn{A2} & The assistant explains the session-scoped setting, gives a cautious cleanup pattern, contrasts \texttt{XACT\_ABORT} with \texttt{TRY}/\texttt{CATCH} and sketches a safer transaction pattern with \texttt{BEGIN TRY}, \texttt{COMMIT}, \texttt{BEGIN CATCH} and \texttt{ROLLBACK}. & \scaffoldbadge\quad \casecode{I2}\ \casecode{M3}\ \casecode{M4}\ \casecode{M5}\par The reply combines \supportkey{procedure}, \supportkey{explanation} and an adaptable model. \\
\bottomrule
\end{tabularx}
\end{minipage}}
\end{adjustbox}
\end{table*}

\begin{table*}[p]
\footnotesize
\centering
\caption{\textbf{Complete de-identified writing case with turn-level labels.} The illustrative transcript is derived from a retained conversation and lightly edited while preserving turn sequence, task structure and annotation decisions. Colored badges encode role, user engagement level and assistant support category.}
\label{tab:complete_writing_case}
\setlength{\fboxsep}{5pt}
\setlength{\fboxrule}{0.6pt}
\begin{adjustbox}{center,max width=0.98\textwidth}
\fcolorbox{caseborder}{casefill}{%
\begin{minipage}{0.955\textwidth}
\setlength{\tabcolsep}{3.5pt}
\renewcommand{\arraystretch}{1.18}
\begingroup
\setlength{\fboxsep}{3pt}
\colorbox{caseheadfill}{%
\begin{minipage}{0.975\linewidth}
\textcolor{caseblue}{\textbf{Writing case}}\quad \textbf{Formal report justification}\\[-1pt]
{\scriptsize WildChat writing; intentional framing; six turns. The exchange moves from direct explanation use to constructive reframing, then to direct drafting under a new style constraint.}
\end{minipage}}
\endgroup
\vspace{5pt}

\caselegend
\vspace{5pt}

\begin{tabularx}{\linewidth}{>{\raggedright\arraybackslash}p{0.115\linewidth}>{\raggedright\arraybackslash}X>{\raggedright\arraybackslash}p{0.305\linewidth}}
\toprule
\textbf{Turn} & \textbf{De-identified transcript} & \textbf{Label and annotation cue} \\
\midrule
\userturn{U1} & ``Help me explain why a machine-learning dataset should include difficult samples and not just clean samples, so that the trained model will be more robust. Give me a more detailed explanation.'' & \activebadge\par The user requests detail but does not yet evaluate, revise or extend an idea. \\
\addlinespace[5pt]
\assistantturn{A1} & The assistant gives a direct explanation: diverse and difficult samples expose the model to real-world variation, reduce dependence on overly clean patterns and improve generalization to noisier cases. & \referencebadge\par The reply supplies requested content without structuring the user's next reasoning or writing process. \\
\addlinespace[5pt]
\userturn{U2} & ``Suppose I am explaining this to someone who does not believe that a model trained only on clean data generalizes worse than one trained with diverse, difficult data. Help me explain why this is the case in more detail.'' & \constructivebadge\par The user \casekey{reframes the task} around persuasion and justification, extending the explanation into a more demanding context. \\
\addlinespace[5pt]
\assistantturn{A2} & The assistant explains the mechanism with an analogy and rationale: clean-only training narrows the problem space, whereas difficult examples expose edge cases, noise and variation that make the learned decision boundary less brittle. & \scaffoldbadge\quad \casecode{I2}\ \casecode{M4}\par The reply teaches the \supportkey{conceptual mechanism} behind robustness rather than merely restating the claim. \\
\addlinespace[5pt]
\userturn{U3} & ``I need a different type of explanation for a formal report. I should justify the need for a diverse dataset without relying on analogies. Can you help me do that?'' & \activebadge\par The user specifies a \casekey{new writing constraint} and asks for a format shift, mainly as directive refinement. \\
\addlinespace[5pt]
\assistantturn{A3} & The assistant provides a formal report-style explanation linking sample diversity to generalization, deployment variation, robustness and evaluation risk. & \referencebadge\par The reply directly supplies formal wording rather than adding process guidance or feedback on the user's own draft. \\
\bottomrule
\end{tabularx}
\end{minipage}}
\end{adjustbox}
\end{table*}
